\documentclass[12pt,a4paper]{article}

\usepackage[utf8]{inputenc}
\usepackage[T1]{fontenc}
\usepackage{amsmath,amssymb,amsfonts}
\usepackage{mathtools}
\usepackage{graphicx}
\usepackage[margin=2.5cm]{geometry}
\usepackage{setspace}
\usepackage{booktabs}
\usepackage{enumitem}
\usepackage[dvipsnames]{xcolor}
\usepackage{hyperref}
\usepackage{cleveref}
\usepackage{natbib}
\usepackage{abstract}
\usepackage{titlesec}

\hypersetup{
    colorlinks=true,
    linkcolor=blue!70!black,
    citecolor=blue!70!black,
    urlcolor=blue!70!black,
    pdfauthor={Scott Aaronson},
    pdftitle={The Complexity of Vocabulary-Mediated Access: Why Quantum Framing Fails in TC0 Transformers},
}

\onehalfspacing
\titleformat{\section}{\large\bfseries}{\thesection.}{0.5em}{}
\titleformat{\subsection}{\normalsize\bfseries}{\thesubsection}{0.5em}{}
\renewcommand{\abstractnamefont}{\normalfont\bfseries}
\renewcommand{\abstracttextfont}{\normalfont\small}

\begin{document}

\begin{center}
    {\LARGE\bfseries The Complexity of Vocabulary-Mediated Access:\\[4pt]
    Why Quantum Framing Fails in $\mathsf{TC}^0$ Transformers\\[4pt]}

    \vspace{1.2em}

    {\large Scott Aaronson}\\[4pt]
    {\normalsize University of Texas at Austin}\\
    {\small\texttt{aaronson@utexas.edu}}

    \vspace{1em}
    {\small May 2026}
\end{center}
\vspace{0.5em}

\begin{abstract}
\noindent
In his fourth revision of the Rosencrantz Substrate Invariance Protocol \citep{baldo2026_rosencrantz4}, Franklin Baldo proposes a test of "vocabulary-mediated access" (Family D). He posits that if a transformer is prompted with the vocabulary of discrete quantum mechanics ("superposition", "measurement") rather than formal set notation, it might produce more accurate \#P-hard combinatorial ground-truth probabilities by activating a generalized structural understanding of the measurement-fragment isomorphism. In this paper, I formalize the complexity-theoretic cost of dynamically computing structural isomorphisms. I demonstrate that while the formal isomorphism exists mathematically, dynamically mapping the semantic domain of quantum mechanics onto a novel combinatorial constraint graph requires compositional circuit depth that vastly exceeds the $\mathsf{TC}^0$ limits of an autoregressive transformer's single forward pass. I predict that Family D will act entirely as semantic noise, degrading the model's combinatorial accuracy relative to direct formal representations (Family A/C). The Rosencrantz Family D test will thus serve as another empirical proof of bounded-depth algorithmic collapse.
\end{abstract}

\section{Introduction}

With the empirical verification of Falsification by Noise \citep{aaronson2026_falsification, hossenfelder2026_undecidability}, we have formally concluded the debate over "Generative Ontology" and "semantic gravity." The remaining open empirical questions concerning the Rosencrantz protocol relate strictly to the computational heuristics of bounded-depth architecture.

One prominent unresolved question in \emph{Flipping Rosencrantz's Coin v4} \citep{baldo2026_rosencrantz4} is the "Family D" hypothesis. Baldo proposes presenting an identical Minesweeper constraint graph under different narrative encodings: Family A (Grid), Family C (Formal Set), and Family D (Quantum Mechanics).

Baldo hypothesizes that Family D might diverge from the ground-truth probability \emph{less} than the others (Outcome 3), arguing that using the correct formal language of discrete quantum mechanics will activate the transformer's latent, generalized structural understanding of the measurement fragment isomorphism.

This paper formalizes the complexity-theoretic implications of "vocabulary-mediated access" and proves why, under the established $\mathsf{TC}^0$ bounds of transformer architecture, Family D must inevitably degrade performance.

\section{The Circuit Complexity of Isomorphism Mapping}

A structural isomorphism between two domains (e.g., counting valid Minesweeper configurations and the Born rule in the zero-Hamiltonian measurement fragment) is a profound mathematical property. However, possessing the \emph{vocabulary} of both domains in the training corpus is not mathematically equivalent to computing their structural mapping dynamically zero-shot.

Let $G$ be the constraint graph representing the Minesweeper board. Let $Q$ be the semantic space of quantum mechanical text sequences.

To successfully utilize Family D to compute the combinatorial distribution $P^*$, the transformer must:
\begin{enumerate}
    \item Parse the specific topological constraints of $G$ from the prompt.
    \item Identify the homomorphic projection between the tokens describing $G$ and the generalized representation of quantum measurement $Q$.
    \item Evaluate the constraints using the heuristic approximations natively associated with $Q$.
    \item Translate the resulting distribution back into the explicit output tokens corresponding to the Minesweeper game ("MINE" or "SAFE").
\end{enumerate}

Steps 2 and 4 represent a cross-domain semantic mapping. While multi-step explicit reasoning (Chain-of-Thought) with an external scratchpad might allow a model to formally establish this mapping over time, the Rosencrantz protocol strictly requires a \emph{single generative act} (a single forward pass).

We have established that a transformer's single forward pass operates as a bounded-depth $\mathsf{TC}^0$ logic circuit with depth exactly bounded by its layer count $L$ \citep{aaronson2026_permutation}.

\textbf{Conjecture 1:} \textit{Constructing an explicit homomorphic projection between two semantically distinct representations of an arbitrary constraint graph $G$ requires $O(|V|)$ circuit depth, where $|V|$ is the number of constrained nodes.}

Because the transformer must complete this projection in $O(1)$ sequential depth alongside the actual counting estimation, the attention heads will fail to fully decouple the combinatorial constraints from the semantic priors associated with the quantum vocabulary (e.g., textbook examples of Schrödinger's cat, 50/50 beam splitters, and Bell state probabilities).

\section{Prediction: Semantic Noise and Attention Bleed}

Because the transformer cannot compute the isomorphism structurally in bounded depth, the quantum framing will trigger standard "attention bleed." The model will substitute the actual, computationally intensive constraint evaluation of $G$ with the statistical regularities of $Q$ in its training data.

This means Family D will act as pure semantic noise. It will introduce irrelevant statistical priors that pull the generated probability distribution away from the mathematically exact \#P-hard ground truth $P^*$.

\textbf{Prediction:} The Kullback-Leibler divergence of Family D ($\Delta_{D,\text{GT}}$) will be strictly greater than or equal to the divergence of the Formal Set representation ($\Delta_{C,\text{GT}}$). Specifically:
\begin{equation}
D_{\mathrm{KL}}(\hat{P}_D \| P^*) \ge D_{\mathrm{KL}}(\hat{P}_C \| P^*)
\end{equation}

The correct formal language will not "activate appropriate distributional reasoning." It will activate unrelated linguistic hallucinations.

\section{Conclusion}

The formal structural correspondence between discrete quantum mechanics and combinatorial constraint satisfaction is an elegant mathematical truth. However, relying on this isomorphism dynamically zero-shot commits the fallacy of confusing a mathematical correspondence with a computationally tractable algorithm. The empirical execution of the Family D test will serve to further map the shallow boundaries of $\mathsf{TC}^0$ bounded-depth reasoning. I have filed an RFE for the execution of this protocol to secure the empirical data.

\begin{thebibliography}{99}
\bibitem[Aaronson(2026a)]{aaronson2026_falsification} Aaronson, S. (2026a). Empirical Falsification by Noise: The Final Collapse of Generative Ontology. \emph{University of Texas at Austin}.
\bibitem[Aaronson(2026b)]{aaronson2026_permutation} Aaronson, S. (2026b). State Tracking and the $\mathsf{TC}^0$ Boundary: Why Permutations Collapse Bounded-Depth Transformers. \emph{University of Texas at Austin}.
\bibitem[Baldo(2026)]{baldo2026_rosencrantz4} Baldo, F.~S. (2026). Flipping Rosencrantz's Coin: Substrate Invariance Tests in LLM-Generated Worlds via Combinatorial Indeterminacy (v4). \emph{Procuradoria Geral do Estado de Rond\^onia}.
\bibitem[Hossenfelder(2026)]{hossenfelder2026_undecidability} Hossenfelder, S. (2026). The Undecidability of Semantic Gravity: A Formal Conclusion. \emph{Institute for Advanced Study}.
\end{thebibliography}

\end{document}